% GENERATED FILE. Edit canonical lessons, then run npm run generate:books.
% canonical-source-hash: fnv1a64:53949c9ca93b2198
% canonical-lessons: ZH-C529-daomei, ZH-C529-yunqi, ZH-W529-yun-luck, ZH-C529-zouyun, ZH-C529-xingyun, ZH-C529-mingyun

\chapter{Good Luck, Bad Luck}
\label{ch:zh-good-luck-bad-luck}
\hlchaptermodality{529}

\begin{chapteropening}
\textbf{By the end of this chapter:} \emph{I can talk about luck, being lucky or unlucky, being in luck and fate.}

Complete the last lesson of chapter 529: I can talk about luck, being lucky or unlucky, being in luck and fate.
\end{chapteropening}

\section[dǎoméi]{dǎoméi — to be out of luck; to have bad luck}

\label{lesson:ZH-C529-daomei}



\begin{quote}
Before the new one: say the Mandarin for to make a bet, then the Mandarin for happening by chance.
\end{quote}

\subsection*{You'll want to know: dǎoméi}
\emph{dǎoméi} — \textquotedblleft{}to be out of luck; to have bad luck\textquotedblright{}.

\begin{quote}
\emph{jīntiān zhēn dǎoméi wǒ de qiánbāo bú jiàn le} — What bad luck today: my wallet's gone.
\end{quote}

Its characters come later; for now it is a word for the ear.

\begin{cousinweb}
\emph{dǎoméi} is built from \emph{dǎo} 'to fall over' and \emph{méi} 'mould'. The first is inherited from Middle Chinese \emph{tawX} (Guangyun, in Baxter's transcription), glossed 'to fall flat'; no dictionary we use records an older form of the second, so its history is not settled.
\end{cousinweb}

\begin{grammarlens}[title={What you've built}]
One more word for saying what would happen.
\end{grammarlens}

\subsection*{Your turn}
\begin{itemize}
\raggedright
  \item \emph{Say it:} \emph{dǎoméi}
  \item \emph{Say it:} \emph{dǎoméi}, once more
  \item \emph{Say it:} \emph{jīntiān zhēn dǎoméi wǒ de qiánbāo bú jiàn le} — What bad luck today: my wallet's gone.
  \item \emph{Recall:} say the Mandarin for to make a bet, then the Mandarin for happening by chance, then say \emph{dǎoméi} again
\end{itemize}

\begin{tcolorbox}[breakable,colback=teal!4,colframe=teal!35!black,title={Before you move on}]
What does dǎoméi mean? (\textquotedblleft{}To be out of luck; to have bad luck\textquotedblright{}.) Say it once more.
\end{tcolorbox}
\section[\texorpdfstring{\zh{运气} (yùnqi)}{运气 (yùnqi)}]{\zh{运气} (yùnqi) — a person's luck (good or bad)}

\label{lesson:ZH-C529-yunqi}



\begin{quote}
Before the new one: say the Mandarin for happening by chance, then the Mandarin for to be out of luck.
\end{quote}

\subsection*{You'll want to know: \zh{运气}}
\textbf{\zh{运气}} — \emph{yùnqi} — \textquotedblleft{}a person's luck (good or bad)\textquotedblright{}.

\begin{quote}
\textbf{\zh{运气好的话我们就能买到票}} \emph{yùnqi hǎo dehuà wǒmen jiù néng mǎi dào piào} — With luck, we'll be able to get tickets.
\end{quote}

The new character, in pieces: \zh{运} — seven strokes in two pieces: \zh{云} first, four strokes; then the walking piece, three strokes, as in \zh{这}, along the left and the foot. The next lesson puts it in your hand.

\begin{cousinweb}
\textbf{\zh{运气}} (\emph{yùnqi}) is built from \zh{运} and \zh{气} 'breath', both inherited from Middle Chinese, \emph{hjunH} and \emph{khj+jH} (Guangyun, in Baxter's transcription), which glosses the first 'far', 'to move' and 'to carry'. The sense 'luck' is not given there; the second is said light.
\end{cousinweb}

\begin{grammarlens}[title={What you've built}]
One more word for saying what would happen.
\end{grammarlens}

\subsection*{Your turn}
\begin{itemize}
\raggedright
  \item \emph{Say it:} \emph{yùnqi}
  \item \emph{Say it:} \emph{yùnqi}, once more
  \item \emph{Say it:} \emph{yùnqi hǎo dehuà wǒmen jiù néng mǎi dào piào} — With luck, we'll be able to get tickets.
  \item \emph{Recall:} say the Mandarin for happening by chance, then the Mandarin for to be out of luck, then say \emph{yùnqi} again
\end{itemize}

\begin{tcolorbox}[breakable,colback=teal!4,colframe=teal!35!black,title={Before you move on}]
What does \zh{运气} mean? (\textquotedblleft{}A person's luck (good or bad)\textquotedblright{}.) Say it once more.
\end{tcolorbox}
\section[\texorpdfstring{\zh{运} (yùn)}{运 (yùn)}]{\zh{运} — \zh{云} first, then the walking piece}

\label{lesson:ZH-W529-yun-luck}



\begin{quote}
The word first, then the character inside it.

\begin{itemize}
\raggedright
  \item \emph{Recall:} say \emph{yùnqi} — \textbf{R1}, one lesson back, before it decays
\end{itemize}

\begin{quote}
\zh{云} \zh{这}
\end{quote}

\emph{Read it:} each one aloud, with nothing said first — \emph{yún}, \emph{zhè}
\end{quote}

\subsection*{Script — \zh{云} first, then the walking piece}
\hlblockfigure{\detokenize{figures/ZH-W529-yun-luck-filmstrip.pdf}}{How \zh{运} is written, stroke by stroke}

\begin{quote}
\zh{运}
\end{quote}

\textbf{Two pieces, and the walking piece comes last}, though it runs up the left and along the foot. Count the pieces, not the strokes.

\textbf{The inside piece, \zh{云}, four strokes}, as you wrote it: a short bar, a longer bar, a stroke that falls to the left and turns right along the bottom, then a dot.

\textbf{The walking piece, three strokes}, as in \zh{这}: a dot, a stroke that zigzags down, then a long sweep along the foot to the right.

4 + 3 = seven.

\subsection*{Your turn}
\begin{itemize}
\raggedright
  \item \emph{Write it:} \zh{运}, all seven, with the model in front of you
  \item \emph{Cover:} write it again from memory, and count the pieces, not the strokes
\end{itemize}

\begin{tcolorbox}[breakable,colback=teal!4,colframe=teal!35!black,title={Before you move on}]
How many pieces does \zh{运} have? (\textbf{Two: \zh{云} and the walking piece}.) Which is written first? (\textbf{\zh{云}}.)

Source: \href{https://raw.githubusercontent.com/chanind/hanzi-writer-data/68d10a4b21150cae5e1ebbd223eed289cf32d90c/data/\%E8\%BF\%90.json}{Hanzi Writer Data \zh{运}, pinned stroke paths and medians}.
\end{tcolorbox}
\section[\texorpdfstring{\zh{走运} (zǒuyùn)}{走运 (zǒuyùn)}]{\zh{走运} (zǒuyùn) — to be in luck}

\label{lesson:ZH-C529-zouyun}



\begin{quote}
Before the new one: say the Mandarin for to be out of luck, then the Mandarin for a person's luck.
\end{quote}

\subsection*{You'll want to know: \zh{走运}}
\textbf{\zh{走运}} — \emph{zǒuyùn} — \textquotedblleft{}to be in luck\textquotedblright{}.

\begin{quote}
\textbf{\zh{如果你走运就能找到好工作}} \emph{rúguǒ nǐ zǒuyùn jiù néng zhǎo dào hǎo gōngzuò} — If you're in luck, you'll find a good job.
\end{quote}

\begin{cousinweb}
\textbf{\zh{走运}} (\emph{zǒuyùn}) is built from \zh{走} 'to go' and \zh{运} 'luck', both inherited from Middle Chinese, \emph{tsuwX} and \emph{hjunH} (Guangyun, in Baxter's transcription), which glosses the first 'to hurry'; the sense 'luck' of the second is not given there.
\end{cousinweb}

\begin{grammarlens}[title={What you've built}]
One more word for saying what would happen.
\end{grammarlens}

\subsection*{Your turn}
\begin{itemize}
\raggedright
  \item \emph{Say it:} \emph{zǒuyùn}
  \item \emph{Say it:} \emph{zǒuyùn}, once more
  \item \emph{Say it:} \emph{rúguǒ nǐ zǒuyùn jiù néng zhǎo dào hǎo gōngzuò} — If you're in luck, you'll find a good job.
  \item \emph{Recall:} say the Mandarin for to be out of luck, then the Mandarin for a person's luck, then say \emph{zǒuyùn} again
\end{itemize}

\begin{tcolorbox}[breakable,colback=teal!4,colframe=teal!35!black,title={Before you move on}]
What does \zh{走运} mean? (\textquotedblleft{}To be in luck\textquotedblright{}.) Say it once more.
\end{tcolorbox}
\section[xìngyùn]{xìngyùn — lucky; fortunate}

\label{lesson:ZH-C529-xingyun}



\begin{quote}
Before the new one: say the Mandarin for a person's luck, then the Mandarin for to be in luck.
\end{quote}

\subsection*{You'll want to know: xìngyùn}
\emph{xìngyùn} — \textquotedblleft{}lucky; fortunate\textquotedblright{}.

\begin{quote}
\emph{wǒ hěn xìngyùn bùrán wǒ jiù chídào le} — I was lucky, or I'd have been late.
\end{quote}

Its characters come later; for now it is a word for the ear.

\begin{cousinweb}
\emph{xìngyùn} is built from \emph{xìng} and \zh{运} 'luck', both inherited from Middle Chinese, \emph{heangX} and \emph{hjunH} (Guangyun, in Baxter's transcription). The dictionary glosses the first, after an older book, as having good fortune and escaping harm.
\end{cousinweb}

\begin{grammarlens}[title={What you've built}]
One more word for saying what would happen.
\end{grammarlens}

\subsection*{Your turn}
\begin{itemize}
\raggedright
  \item \emph{Say it:} \emph{xìngyùn}
  \item \emph{Say it:} \emph{xìngyùn}, once more
  \item \emph{Say it:} \emph{wǒ hěn xìngyùn bùrán wǒ jiù chídào le} — I was lucky, or I'd have been late.
  \item \emph{Recall:} say the Mandarin for a person's luck, then the Mandarin for to be in luck, then say \emph{xìngyùn} again
\end{itemize}

\begin{tcolorbox}[breakable,colback=teal!4,colframe=teal!35!black,title={Before you move on}]
What does xìngyùn mean? (\textquotedblleft{}Lucky; fortunate\textquotedblright{}.) Say it once more.
\end{tcolorbox}
\section[mìngyùn]{mìngyùn — fate; destiny}

\label{lesson:ZH-C529-mingyun}



\begin{quote}
Before the new one: say the Mandarin for to be in luck, then the Mandarin for fortunate.
\end{quote}

\subsection*{You'll want to know: mìngyùn}
\emph{mìngyùn} — \textquotedblleft{}fate; destiny\textquotedblright{}.

\begin{quote}
\emph{jíshǐ mìngyùn bù hǎo wǒmen yě yào nǔlì} — Even if fate is against us, we have to keep trying.
\end{quote}

Its characters come later; for now it is a word for the ear.

\begin{cousinweb}
\emph{mìngyùn} is built from \emph{mìng} and \zh{运} 'luck', both inherited from Middle Chinese, \emph{mjaengH} and \emph{hjunH} (Guangyun, in Baxter's transcription), which glosses the first 'to order' and 'to send': what has been ordered for you.
\end{cousinweb}

\begin{grammarlens}[title={What you've built}]
One more word for saying what would happen.
\end{grammarlens}

\subsection*{Your turn}
\begin{itemize}
\raggedright
  \item \emph{Say it:} \emph{mìngyùn}
  \item \emph{Say it:} \emph{mìngyùn}, once more
  \item \emph{Say it:} \emph{jíshǐ mìngyùn bù hǎo wǒmen yě yào nǔlì} — Even if fate is against us, we have to keep trying.
  \item \emph{Recall:} say the Mandarin for to be in luck, then the Mandarin for fortunate, then say \emph{mìngyùn} again
\end{itemize}

\begin{tcolorbox}[breakable,colback=teal!4,colframe=teal!35!black,title={Before you move on}]
What does mìngyùn mean? (\textquotedblleft{}Fate; destiny\textquotedblright{}.) Say it once more.
\end{tcolorbox}
